\section{Exact synthesis into the paper gate set}
\label{sec:exact}

This section is the first half of the universality proof: every $n$-qubit
unitary, $n \ge 1$, is \emph{exactly} a circuit over
$\{C(X), H, S, S^\dagger, R_z(\theta)\}$ up to global phase.  It follows
\texttt{universal\_new\_gates.tex}, Lemma
\texttt{clifford-plus-rx-is-universal}: an induction on the number of qubits
whose base case is one qubit.  Every node in this section is proved.

Throughout, $m$ control wires and one target wire carry two families of
\emph{uniformly controlled} rotations.  Writing $r$ for the control register
value, $\operatorname{controlledRzFamily} m\, \alpha$ applies
$R_z(\alpha_r)$ to the target, and $\operatorname{controlledRyFamily} m\, \theta$
applies $R_y(\theta_r)$; and $\operatorname{firstQubitBlockDiag} m\, U_0\, U_1$
is the gate applying $U_0$ or $U_1$ to the lower $m$ wires according to the top
wire.

\subsection{The recursion step}

\begin{lemma}[Cosine--sine decomposition, Paige--Wei]
  \label{lem:cs_step}
  \lean{TwoControl.Clifford.Universal.general_cosine_sine_step_exists}
  \leanok
  Let $n \ge 1$ and let $U$ be a $2^n \times 2^n$ unitary.  Then there are
  unitaries $P, R, Q$ with $U = P R Q$, where $P$ and $Q$ are block diagonal on
  the first qubit and $R$ is a uniformly controlled $R_y$ family:
  \[
    P = \begin{pmatrix} P_0 & 0 \\ 0 & P_1\end{pmatrix},\qquad
    Q = \begin{pmatrix} Q_0 & 0 \\ 0 & Q_1\end{pmatrix},\qquad
    R = \begin{pmatrix} C & -S \\ S & C \end{pmatrix},
  \]
  with $C$ and $S$ the diagonal matrices of $\cos(\theta_r/2)$ and
  $\sin(\theta_r/2)$.
\end{lemma}

\begin{proof}
  \leanok
  Split $U$ into $2 \times 2$ blocks of size $2^{n-1}$.  Unitarity makes the
  four blocks simultaneously reducible: diagonalizing $A^\dagger A$ produces a
  common cosine/sine pair, and the corresponding unitary factors assemble into
  the block-diagonal outer terms.  The construction is carried out directly on
  blocks rather than imported, so the general case is available for every
  $n \ge 1$; the two-qubit instance is also recorded separately as
  Lemma~\ref{lem:cs_step_2q} and is what the $n = 2$ layer of the recursion
  uses.
\end{proof}

\begin{lemma}
  \label{lem:cs_step_2q}
  \lean{TwoControl.Clifford.Universal.two_qubit_cosine_sine_step}
  \leanok
  The two-qubit instance of Lemma~\ref{lem:cs_step}, stated with an explicit
  witness structure.
\end{lemma}

\begin{proof}
  \leanok
  \uses{lem:cs_exists_2q}
  This is the cosine--sine decomposition of Lemma~\ref{lem:cs_exists_2q},
  repackaged in the $n$-qubit witness shape.
\end{proof}

\begin{lemma}[$R_y$ via $R_z$]
  \label{lem:ry_via_rz}
  \lean{TwoControl.Clifford.lemma3_ry_via_rz}
  \leanok
  For every real $\theta$,
  \[
    R_y(\theta) = S^\dagger\, H\, R_z(-\theta)\, H\, S.
  \]
\end{lemma}

\begin{proof}
  \leanok
  Expand both sides as $2 \times 2$ matrices.  The identities $HZH = X$ and
  $S^\dagger X S = -Y$ reduce the right-hand side to the defining matrix of
  $R_y(\theta)$.
\end{proof}

\begin{lemma}
  \label{lem:cry_via_crz}
  \lean{TwoControl.Clifford.Universal.controlled_ry_family_via_controlled_rz}
  \leanok
  A uniformly controlled $R_y$ family is a uniformly controlled $R_z$ family
  conjugated by local Clifford gates on the target wire:
  \[
    \operatorname{controlledRyFamily} m\, \theta
      = (S^\dagger \otimes I)(H \otimes I)\,
        \operatorname{controlledRzFamily} m\, (-\theta)\,
        (H \otimes I)(S \otimes I).
  \]
\end{lemma}

\begin{proof}
  \leanok
  \uses{lem:ry_via_rz}
  Apply Lemma~\ref{lem:ry_via_rz} entrywise in the control index: both sides
  are block matrices whose blocks are diagonal, and the corresponding diagonal
  entries agree by the one-qubit identity.
\end{proof}

\begin{lemma}[Demultiplexing, Shende--Bullock--Markov]
  \label{lem:demux}
  \lean{TwoControl.Clifford.Universal.general_demultiplexing_step}
  \leanok
  Let $m \ge 1$ and let $U_0, U_1$ be $2^m \times 2^m$ unitaries.  Then there
  are unitaries $P, Q$ on $m$ qubits and angles $\alpha$ with
  \[
    \operatorname{firstQubitBlockDiag} m\, U_0\, U_1
      = (I \otimes Q)\,
        \operatorname{controlledRzFamily} m\, \alpha\,
        (I \otimes P).
  \]
\end{lemma}

\begin{proof}
  \leanok
  Diagonalize $W = U_1 U_0^\dagger = Q\,\operatorname{diag}(z)\,Q^\dagger$ and
  set $\alpha_i = \arg z_i$.  Taking $D_{0}, D_1$ to be the diagonal matrices
  of $e^{\mp i\alpha_i/2}$ and $P = D_0^\dagger Q^\dagger U_0$ gives the stated
  factorization; the middle factor is by construction a uniformly controlled
  $R_z$ family.
\end{proof}

\begin{lemma}[Controlled-$R_z$ recursion, M\"ott\"onen et al.]
  \label{lem:mottonen}
  \lean{TwoControl.Clifford.Universal.controlled_rz_reduction_step}
  \leanok
  A uniformly controlled $R_z$ family with $m+1$ controls splits into two
  families with $m$ controls and two $C(X)$ gates: there exist $\beta, \gamma$
  and an embedded $C(X)$ with
  \[
    \operatorname{controlledRzFamily}(m{+}1)\,\alpha
      = C(X)\cdot \operatorname{controlledRzFamily} m\, \beta
        \cdot C(X)\cdot \operatorname{controlledRzFamily} m\, \gamma
  \]
  after lifting the smaller families onto the appropriate wires.
\end{lemma}

\begin{proof}
  \leanok
  Take $\beta_r$ and $\gamma_r$ to be the half-difference and half-sum of
  $\alpha$ at the two extensions of the control index $r$.  Both sides are
  block diagonal over $r$, and on each block the identity is the standard
  two-$C(X)$ realization of a controlled $R_z$ pair.
\end{proof}

\begin{lemma}
  \label{lem:crz_family_synth}
  \lean{TwoControl.Clifford.Universal.synthesizes_controlled_rz_family_cliffordRz}
  \leanok
  Every uniformly controlled $R_z$ family is \emph{exactly} a circuit over the
  paper gate set.
\end{lemma}

\begin{proof}
  \leanok
  \uses{lem:mottonen, def:gate_cliffordrz, def:circuit_over}
  Induct on the number of controls using Lemma~\ref{lem:mottonen}.  The base
  case is a single $R_z(\theta)$ on one wire, which is a paper gate; each
  recursion step contributes two embedded $C(X)$ gates.
\end{proof}

\begin{lemma}
  \label{lem:cry_family_synth}
  \lean{TwoControl.Clifford.Universal.synthesizes_controlled_ry_family_cliffordRz}
  \leanok
  Every uniformly controlled $R_y$ family is exactly a circuit over the paper
  gate set.
\end{lemma}

\begin{proof}
  \leanok
  \uses{lem:cry_via_crz, lem:crz_family_synth}
  Rewrite with Lemma~\ref{lem:cry_via_crz} and apply
  Lemma~\ref{lem:crz_family_synth}; the four conjugating factors $S^\dagger$,
  $H$, $H$, $S$ are themselves paper gates.
\end{proof}

\subsection{The one-qubit base case}

\begin{lemma}[$ZYZ$ decomposition]
  \label{lem:oneq_euler}
  \lean{TwoControl.Clifford.one_qubit_euler_rz_ry_rz_up_to_global_phase}
  \leanok
  For every one-qubit unitary $U$ there are real $\alpha, \beta, \gamma$ and a
  scalar $z$ with $\lvert z\rvert = 1$ such that
  $U = z\,R_z(\alpha)\,R_y(\beta)\,R_z(\gamma)$.
\end{lemma}

\begin{proof}
  \leanok
  This is the Nielsen--Chuang $ZYZ$ decomposition.  Normalize $U$ to determinant
  one by extracting a scalar, and read off the three Euler angles from the
  entries of the resulting special unitary.
\end{proof}

\begin{lemma}[Paper base case]
  \label{lem:oneq_base}
  \lean{TwoControl.Clifford.one_qubit_exact_clifford_rz}
  \leanok
  For every one-qubit unitary $U$ there are real $\alpha, \beta, \gamma$ with
  \[
    U \;\sim\; R_z(\alpha)\, S^\dagger H\, R_z(-\beta)\, H S\, R_z(\gamma),
  \]
  a circuit over the paper gate set of length seven.
\end{lemma}

\begin{proof}
  \leanok
  \uses{lem:oneq_euler, lem:ry_via_rz, def:global_phase}
  Apply Lemma~\ref{lem:oneq_euler} and replace the middle $R_y(\beta)$ using
  Lemma~\ref{lem:ry_via_rz}.  The leftover scalar is exactly the global phase.
\end{proof}

\subsection{The induction}

\begin{theorem}[Exact synthesis over $\{C(X),H,S,S^\dagger,R_z\}$]
  \label{thm:clifford_rz_universal}
  \lean{TwoControl.Clifford.Universal.clifford_rz_universal}
  \leanok
  For every $n \ge 1$ and every $2^n \times 2^n$ unitary $U$, some circuit over
  the paper gate set has matrix globally phase equivalent to $U$.
\end{theorem}

\begin{proof}
  \leanok
  \uses{lem:oneq_base, lem:cs_step, lem:demux, lem:crz_family_synth,
        lem:cry_family_synth, def:synthesizes_phase}
  Strong induction on $n$.  For $n = 1$ this is Lemma~\ref{lem:oneq_base}.

  For $n = m+1$ with $m \ge 1$, apply Lemma~\ref{lem:cs_step} to write
  $U = PRQ$ with $P, Q$ block diagonal on the first qubit and $R$ a uniformly
  controlled $R_y$ family.  The middle factor is handled by
  Lemma~\ref{lem:cry_family_synth}.  For each outer factor, unitarity of the
  block-diagonal matrix gives unitarity of its two blocks; apply
  Lemma~\ref{lem:demux} to turn the factor into two $m$-qubit gates
  surrounding a uniformly controlled $R_z$ family, synthesize the $m$-qubit
  gates by the induction hypothesis, and the middle family by
  Lemma~\ref{lem:crz_family_synth}.  Concatenating the three circuits and
  multiplying the accumulated scalars gives a single global phase.

  Wire plumbing inside the recursion moves gates between the top and lower
  registers; the required $\operatorname{SWAP}$ gates are expanded as
  $\operatorname{SWAP} = C(X)\,(H\otimes H)\,C(X)\,(H\otimes H)\,C(X)$, which
  stays inside the paper gate set.
\end{proof}

\subsection{Rewriting the Clifford phases as powers of T}

\begin{proposition}
  \label{prop:embedded_s}
  \lean{TwoControl.Clifford.Universal.embedded_phaseS_is_clifford_t}
  \leanok
  Every embedded $S$ gate is synthesized exactly by an embedded Clifford+T
  circuit.
\end{proposition}

\begin{proof}
  \leanok
  \uses{def:gate_cliffordt, def:embedded_one}
  Use the exact identity $S = T^2$ and lift the two-letter one-qubit circuit to
  the selected wire.
\end{proof}

\begin{proposition}
  \label{prop:embedded_sdag}
  \lean{TwoControl.Clifford.Universal.embedded_phaseSdagger_is_clifford_t}
  \leanok
  Every embedded $S^\dagger$ gate is synthesized exactly by an embedded
  Clifford+T circuit.
\end{proposition}

\begin{proof}
  \leanok
  \uses{def:gate_cliffordt, def:embedded_one}
  Use $S^\dagger = T^6$ and lift as before.
\end{proof}

\begin{theorem}
  \label{thm:clifford_rz_to_t_rz}
  \lean{TwoControl.Clifford.Universal.clifford_rz_to_clifford_t_rz}
  \leanok
  If $U$ is synthesized up to global phase over the paper gate set, then it is
  synthesized up to global phase over the intermediate gate set
  $\{C(X), H, T, R_z(\theta)\}$.
\end{theorem}

\begin{proof}
  \leanok
  \uses{prop:embedded_s, prop:embedded_sdag, def:gate_cliffordtrz,
        thm:clifford_rz_universal}
  Replace each factor of the circuit: $C(X)$, $H$, and $R_z(\theta)$ are
  already in the target set, and the two Clifford phases are expanded by
  Propositions~\ref{prop:embedded_s} and~\ref{prop:embedded_sdag}.  Since each
  replacement is exact, the global phase is unchanged.
\end{proof}
